\documentclass[12pt]{article}

\usepackage[margin=1in]{geometry}
\usepackage{fontspec}
\usepackage{hyperref}
\usepackage{xcolor}
\usepackage{setspace}
\usepackage{parskip}
\usepackage{titlesec}

\setmainfont{TeX Gyre Termes}
\hypersetup{
  colorlinks=true,
  linkcolor=black,
  urlcolor=blue,
  citecolor=black
}

\titleformat{\section}{\large\bfseries}{}{0pt}{}
\setstretch{1.15}

\title{Your LaTeX Blog Title}
\author{Your Name}
\date{\today}

\begin{document}

\maketitle

\begin{center}
\textit{A minimal LaTeX blog template for GitHub Pages.}
\end{center}

\section*{Introduction}

This repository is a lightweight starting point for publishing a blog post, essay,
or technical note written in \LaTeX.

You only need to edit \texttt{main.tex}, push to GitHub, and let GitHub Actions
compile and deploy the result.

\section*{Why this template exists}

Many personal sites want the simplicity of writing in \LaTeX, but also want a
clean web page for reading and sharing. This template keeps the workflow small:

\begin{itemize}
  \item write in \LaTeX
  \item compile automatically in CI
  \item publish the PDF with a simple HTML wrapper
\end{itemize}

\section*{How to customize}

Replace this content with your own post. You can freely add sections, equations,
figures, references, or bilingual text.

For example, you can add a link like this:
\href{https://github.com/}{GitHub}.

\section*{Next step}

Turn this into your own blog post and publish it.

\end{document}
